% Generated draft: requires mathematical and source review.
\documentclass[11pt]{article}
\usepackage[margin=25mm]{geometry}
\usepackage{fontspec}
\setmainfont{Latin Modern Roman}
\usepackage{amsmath,amssymb,mathtools,unicode-math}
\setmathfont{Latin Modern Math}
\usepackage{enumitem,needspace,xurl,hyperref}
\setlength{\parindent}{0pt}
\setlength{\parskip}{5pt}
\setlength{\emergencystretch}{3em}
\newcommand{\N}{\mathbb{N}}
\newcommand{\Z}{\mathbb{Z}}
\newcommand{\Q}{\mathbb{Q}}
\newcommand{\R}{\mathbb{R}}
\newcommand{\C}{\mathbb{C}}
\newcommand{\F}{\mathbb{F}}
\newcommand{\A}{\mathbb{A}}
\newcommand{\PP}{\mathbb{P}}
\newcommand{\E}{\mathbb{E}}
\newcommand{\eps}{\varepsilon}
\newcommand{\dd}{\,\mathrm{d}}
\begin{document}
\section*{UnsolvedMath ID: 700953}
\textbf{Source problem code:} catalogue entry
\subsection{Definitions and mathematical statement}
Let \(p\) be a prime. A pro-\(p\) group is a topological group isomorphic to an inverse limit of finite \(p\)-groups. For such a group \(G\), set \(G^{(0)}=G\) and
\[
G^{(n+1)}=\overline{[G^{(n)},G^{(n)}]}\qquad(n\geq 0).
\]
Here \(G\) is called soluble if \(G^{(n)}=\{1\}\) for some finite \(n\), and non-soluble otherwise. This is solubility of finite derived length, not merely being an inverse limit of finite soluble groups.

Interpret cohomological dimension as continuous \(p\)-cohomological dimension. Thus \(\operatorname{cd}_p(G)=2\) means that
\[
H^q_{\mathrm{cts}}(G,M)=0
\]
for every discrete \(p\)-primary torsion \(G\)-module \(M\) with continuous action and every integer \(q>2\), while \(H^2_{\mathrm{cts}}(G,M)\neq 0\) for at least one such module. Let \(F_2\) be the abstract free group on two generators, and let \(\widehat{F_2}^{\,(p)}\) be its pro-\(p\) completion, the free pro-\(p\) group of rank two.

The question is whether, for every prime \(p\) and every non-soluble pro-\(p\) group \(G\) with \(\operatorname{cd}_p(G)=2\), there exists a closed subgroup \(H\leq G\) such that
\[
H\cong_{\mathrm{top}}\widehat{F_2}^{\,(p)}.
\]
Equivalently, does every such \(G\) contain a closed non-abelian free pro-\(p\) subgroup? No finite generation hypothesis on \(G\) is supplied. The source does not explicitly specify the cohomological or subgroup-topology conventions; the statement above uses continuous cohomology and closed subgroups with their induced topology, as customary for pro-\(p\) groups.
\subsection{Short English statement}
For every prime \(p\), must each non-soluble pro-\(p\) group of continuous cohomological dimension exactly \(2\) contain a closed non-abelian free pro-\(p\) subgroup? This is a statement-only draft of the question attributed to O. V. Mel'nikov.
\subsection{Sources}
\begin{itemize}
\item Local UnsolvedMath registry, permanent ID \(700953\), exact source problem code \texttt{catalogue entry}, title ``Kourovka Notebook Problem 13.35.'' Supplied upstream revision: \texttt{d371ee917ac9647fbb06cd6fc529deec6f3dda82}.
\item The supplied record attributes the problem to O. V. Mel'nikov, proposed in 1995, and cites \emph{The Kourovka Notebook (Unsolved Problems in Group Theory)}, 13th issue (1995), Problem 13.35.
\item The supplied background also cites the 21st issue edition (2026), identified as \texttt{arXiv:1401.0300v48}. The registry describes the collection as Issues 1--21 (1965--2026), edited by E. I. Khukhro and V. D. Mazurov. No external URL or page reference is supplied.
\item The registry labels this problem ``open'' and records an update date of 2026-10-16. This is supplied registry metadata, not an independently verified assessment of the literature or mathematical status.
\end{itemize}
\end{document}
